%-------------------------------------------------------------------------------
%	SECTION TITLE
%-------------------------------------------------------------------------------
\cvsection{Publications}


%\textit{\footnotesize My publications focus on dualities in logic, especially in connection with non-classical logics.}

\newcounter{pubscounter}
\newcommand\increaseandshowpubs{\addtocounter{pubscounter}{1}\thepubscounter}

%-------------------------------------------------------------------------------
%	CONTENT
%-------------------------------------------------------------------------------



%---------------------------------------------------------
\cvsubsection{\stress{Publications in international peer-reviewed journals}}
%---------------------------------------------------------




\begin{cvpubs}



	
	\cvpub{
		\increaseandshowpubs. M.\ Abbadini, L.\ Spada.\\
		 The unification type of {\L}ukasiewicz logic with a bounded number of variables.\\
		 \textbf{\em The Journal of Symbolic Logic}.
		\gray In press. 
		DOI: \href{https://doi.org/10.1017/jsl.2025.10149}{10.1017/jsl.2025.10149}.
		%Preprint at \href{https://arxiv.org/abs/2504.19011}{arXiv:2504.19011}.
	}
	
	

	\cvpub{
		\increaseandshowpubs. M.\ Abbadini, G.\ Bezhanishvili, L.\ Carai.\\
		On the lack of colimits in various categories arising in pointfree topology and algebraic logic.\\
		\textbf{\em Theory and Applications of Categories}.
		 \gray (2026) 45(20):759--778.
		URL: \url{http://www.tac.mta.ca/tac/volumes/45/20/45-20abs.html}.
	}

	
	
	\cvpub{
		\increaseandshowpubs. M.\ Abbadini, V.\ Marra, L.\ Spada.\\
		 Stone--Gelfand duality for metrically complete lattice-ordered groups.\\
		 \textbf{\em Advances in Mathematics}.
		\gray (2025) 461:110067. 
		DOI: \href{https://doi.org/10.1016/j.aim.2024.110067}{10.1016/j.aim.2024.110067}.
	}



	\cvpub{
		\increaseandshowpubs. M.\ Abbadini, I.\ Di Liberti.\\
		 Duality for coalgebras for Vietoris and monadicity.\\
		 \textbf{\em The Journal of Symbolic Logic}.
		\gray (2025) 90(3):1022--1055.
		 DOI: \href{https://doi.org/10.1017/jsl.2024.14}{10.1017/jsl.2024.14}.
	}

	\cvpub{
		\increaseandshowpubs. M.\ Abbadini, P.\ Aglian\`o, S.\ Fioravanti.\\
		 Varieties of MV-monoids and positive MV-algebras.\\
		 \textbf{\em Journal of Algebra}.
		\gray (2025) 677:690--744. DOI: \href{https://doi.org/10.1016/j.jalgebra.2025.04.027}{10.1016/j.jalgebra.2025.04.027}.
	}
	

	\cvpub{
		\increaseandshowpubs. M.\ Abbadini, F.\ Guffanti.\\
		 Quantifier-free formulas and quantifier alternation depth in doctrines.\\
		 \textbf{\em Journal of Pure and Applied Algebra}.
		\gray (2025) 229(8):108004. DOI: \href{https://doi.org/10.1016/j.jpaa.2025.108004}{10.1016/j.jpaa.2025.108004}.
	}


	\cvpub{
		\increaseandshowpubs. M.\ Abbadini, D.\ Hofmann.\\
		 Barr-coexactness for metric compact Hausdorff spaces.\\
		 \textbf{\em Theory and Applications of Categories}.
		\gray (2025) 44(6):196--226.  URL: \url{http://www.tac.mta.ca/tac/volumes/44/6/44-06abs.html}.%{http://www.tac.mta.ca/tac/volumes/44/6/44-06abs.html}.
		}

	
		
	\cvpub{
		\increaseandshowpubs. M.\ Abbadini, G.\ Bezhanishvili, L.\ Carai.\\
		 Vietoris endofunctor for closed relations and its de Vries dual.\\
		 \textbf{\em Topology Proceedings}.
		\gray (2024) 64:213--250. URL: \url{http://topology.nipissingu.ca/tp/reprints/v64}.%{https://topology.nipissingu.ca/tp/reprints/v64/}. %Preprint at \href{https://arxiv.org/abs/2308.16823}{arXiv:2308.16823}.
	}


	
	
	
	\cvpub{
		\increaseandshowpubs. M.\ Abbadini, G.\ Bezhanishvili, L.\ Carai.\\
		 MacNeille completions of subordination algebras.\\
		 \textbf{\em Cahiers de Topologie et Géométrie Différentielle Catégoriques}.
		\gray (2024) 65(2):151--199. 
		 URL: \url{https://cahierstgdc.com/wp-content/uploads/2024/04/ABBADINI-Al-LXV-2.pdf}.%{https://cahierstgdc.com/{\newline}wp-content/uploads/2024/04/ABBADINI-Al-LXV-2.pdf}.
	}
	
	\cvpub{
		\increaseandshowpubs. M.\ Abbadini, G.\ Bezhanishvili, L.\ Carai.\\
		 A generalization of de Vries duality to closed relations between compact Hausdorff spaces.\\
		 \textbf{\em Topology and its Applications}.
		 \gray (2023) 337:108641.
		 DOI: \href{https://doi.org/10.1016/j.topol.2023.108641}{10.1016/j.topol.2023.108641}.
	}
	
	
	\cvpub{
		\increaseandshowpubs. M.\ Abbadini, L.\ Reggio.\\
		 Barr-exact categories and soft sheaf representations.\\
		 {\em \textbf{Journal of Pure and Applied Algebra}}.
		 \gray 
		(2023) 227(12):107413.
		 DOI: \href{https://doi.org/10.1016/j.jpaa.2023.107413}{10.1016/j.jpaa.2023.107413}.
	}
	
	\cvpub{ \label{i:posMV}
		\increaseandshowpubs. M.\ Abbadini, P.\ Jipsen, T.\ Kroupa, S.\ Vannucci.\\
		 A finite axiomatization of positive MV-algebras.\\
		 {\em \textbf{Algebra Universalis}}.
		 \gray (2022) 83:28.
		 DOI: \href{https://doi.org/10.1007/s00012-022-00776-3}{10.1007/s00012-022-00776-3}.
		%		B\\
		%		\bodyfontlight
		%		\begin{itemize}
			%			\item 
			%			I consider this to be one of my best two papers.\\
			%			It is in the research area of the algebraic study of {\L}ukasiewicz many-valued logic.
			%			An interpretation of {\L}ukasiewicz logic can be given in the framework of Rényi-Ulam games, a variant of the game of Twenty Questions where a fixed amount of lies are allowed in the answers. 
			%			In the traditional game without lies, someone thinks of an object and another person should guess it using twenty yes-or-no questions.
			%			The game with lies is described in Ulam's book ``Adventures of a Mathematician'' and by Rényi’s as a problem of fault-tolerant adaptive search with errors.
			%			Ulam–Rényi games are part of Berlekamp’s communication theory with feedback.
			%			Just as the states of knowledge in the game without lies form a boolean algebra, the states of knowledge in every game with lies form a so-called MV-algebra.
			%			In fact, just like Boolean algebras are the algebras of classical propositional logic, MV-algebras are the algebras of {\L}ukasiewicz many-valued logic, which, in the most general formulation, has the unit interval $[0,1]$ instead of $\{0,1\}$ as the set of truth values.
			%			MV-algebras are then a generalization of Boolean algebras where $\{0,1\}$ is replaced by $[0,1]$.
			%			\\
			%			Algebraic structures of interest in logic are partially ordered.
			%			We can then distinguish the operations that preserve the order (such as unions and intersections) and the operation that don't (such as negation, difference, implication).
			%			This gives interest in the \emph{positive subreducts}, i.e.\ the subsets closed under order-preserving operations.
			%			For example, the positive subreducts of Boolean algebras are precisely the much studied bounded distributive lattices.\\
			%			In this work, we are interested in the positive subreducts of MV-algebras.
			%			In a sense, these are the many-valued version of bounded distributive lattices: Boolean algebras are to MV-algebras as bounded distributive lattices are to positive subreducts of MV-algebras.\\
			%			Our main result is a description of the class of positive subreducts of MV-algebras via finitely many first-order axioms; even more, these axioms are all equational except for one that is an implication of equations.
			%			This result is the best possible in the sense that it cannot be improved to an equational axiomatization.
			%			We rely, among other things, to a result I obtained in the paper ``Equivalence à la Mundici for commutative lattice-ordered monoids'', which contains some of the hard computations.
			%			Our main result solves an open problem in a conference abstract by L.~Cabrer, P.~Jipsen and T.~Kroupa.
			%			I consider it one of my best papers because (i) we obtained the best possible axiomatization for the class of interest, solving a recent open problem, and (ii) the proof is short and elementary.
			%%Content: 
			%%			We solve the open problem in ``L.~M.~Cabrer, P.~Jipsen, T.~Kroupa. Positive
			%%			subreducts in finitely generated varieties of MV-algebras. Presented at SYSMICS,
			%%			Amsterdam, 2019''.
			%%			Just like Boolean algebras provide the algebraic semantics of classical propositional logic, MV-algebras provide the algebraic semantics of {\L}ukasiewicz logic.
			%%			The relevance of the order in the algebraic structures of interest in logic gives interest to the study of order-preserving terms, with the consequent notion of \emph{positive subreducts}.
			%%%			: for example, the positive subreducts of Boolean algebras are precisely bounded distributive lattices.
			%%			Our main result is a finite quasi-equational axiomatization for the class of positive subreducts of MV-algebras.
			%%			It should be noted that this result is the best possible in the sense that it cannot be improved to an equational axiomatization.
			%%%			While the class of positive subreducts of Boolean algebras admits an equational axiomatization (namely, the axioms of bounded distributive lattices), the class of positive subreducts of MV-algebras does not.
			%			\item Review (MathSciNet): \href{https://mathscinet.ams.org/mathscinet/pdf/4447154.pdf}{https://mathscinet.ams.org/mathscinet/pdf/4447154.pdf}.
			%			%
			%		\end{itemize}
	}
	
	
	\cvpub{
		\increaseandshowpubs. M.\ Abbadini, L.\ Spada.\\
		 Are locally finite MV-algebras a variety? \\
		 {\em \textbf{Journal of Pure and Applied Algebra}}.
		 \gray (2022) 226(4):106858.
		 DOI: \href{https://doi.org/10.1016/j.jpaa.2021.106858}{10.1016/j.jpaa.2021.106858}.
		%		A\\
		%		\bodyfontlight
		%		\begin{itemize}
			%%			\item Content: %: for example, they have free algebras on all sets and they are closed under products, subalgebras and homomorphic images.
			%%			We solve an open problem in the book ``D.~Mundici. \emph{Advanced {\L}ukasiewicz calculus.} Trends in Logic Vol.~35. Springer 2011''): \emph{Is the category of locally finite MV-algebras equivalent to an equational class?} (Problem 3, p.\ 235.)
			%%%			The equational classes of algebras, also known as varieties of algebras, are classes of algebras with good properties.
			%%			The answer depends on the numbers of sorts and arities allowed (finite vs infinite).
			%%			Our result excludes a finite equational treatment of locally finite MV-algebras and is a starting point for an infinite equational approach.
			%			\item
			%			Review (MathSciNet): \href{https://mathscinet.ams.org/mathscinet/pdf/4297845.pdf}{https://mathscinet.ams.org/mathscinet/pdf/4297845.pdf}.\\
			%			Review (Zentral Blatt): \href{https://zbmath.org/?q=an\%3A7428823}{https://zbmath.org/?q=an\%3A7428823}.
			%		\end{itemize}
	}
	
	\cvpub{
		\increaseandshowpubs. M.\ Abbadini.\\
		 Equivalence à la Mundici for commutative lattice-ordered monoids.\\
		 {\em \textbf{Algebra Universalis}}.
		 \gray
		 (2021) 82:45.
		 DOI: \href{https://doi.org/10.1007/s00012-021-00736-3}{10.1007/s00012-021-00736-3}.
		%    		B\\
		%		\bodyfontlight
		%		\begin{itemize}
			%%			\item Content: %MV-algebras are the algebraic semantics of {\L}ukasiewicz logic.
			%%			A cornerstone result in the theory of MV-algebras is Mundici's categorical equivalence between unital abelian lattice-ordered groups and MV-algebras (1986).
			%%			We generalize it to a negation-free framework by building a categorical equivalence between unital commutative lattice-ordered monoids and the algebras introduced in the article under the name of \emph{MV-monoidal algebras}.
			%%			This equivalence allows the study of positive subreducts of MV-algebras via the theory of lattice-ordered monoids.
			%%			In fact, jointly with P.~Jipsen, T.~Kroupa, S.~Vannucci, we derived from this equivalence a finite axiomatization of the class of positive subreducts of MV-algebras (M.~Abbadini, P.~Jipsen, T.~Kroupa, S.~Vannucci. A finite axiomatization of positive MV-algebras. {\em Algebra Universalis}, 83:28, 2022).
			%			\item
			%			Review (MathSciNet): \href{https://mathscinet.ams.org/mathscinet/pdf/4280475.pdf}{https://mathscinet.ams.org/mathscinet/pdf/4280475.pdf}.
			%		\end{itemize}
	}
	
	
	\cvpub{
		\increaseandshowpubs. M.\ Abbadini, L.\ Reggio.\\
		 On the axiomatisability of the dual of compact ordered spaces.\\
		 {\em \textbf{Applied Categorical Structures}}.
		 \gray (2020) 28(6):921--934.
		 DOI: \href{https://doi.org/10.1007/s10485-020-09604-y}{10.1007/s10485-020-09604-y}.
		%		B\\
		%		\bodyfontlight
		%		\begin{itemize}
			%			\item 
			%			In this paper, that I consider to be one of the two most representative works of mine, I investigate the duality theory of compact ordered spaces. I build on the work of M.~H.~Stone, who in 1936 established a duality between Boolean algebras and Stone spaces, which are compact Hausdorff spaces with a basis of closed open sets. This duality, known as Stone's Representation Theorem, has had a wide-reaching impact on various areas of mathematics and computer science.\\			
			%%			In the paper, I show that the duality of Priestley spaces can be extended to the category of compact ordered spaces while retaining its algebraic nature. 
			%%			In the paper “The dual of compact ordered spaces is a variety” I had already done this by providing an explicit equational axiomatization of the dual category, based on an algebraic language related to {\L}ukasiewicz many-valued logic. In this paper with Reggio, we provide a new proof. The proof is purely categorical; thus, it does not depend on a particular choice of the dual language. The proof copes with a difficulty that in the unordered case considered by Duskin is not present (i.e.\  the dual of compact ordered spaces is not Mal’cev), requiring new techniques.				
			%%			I provide an explicit equational axiomatization of the dual category, which is based on an algebraic language related to {\L}ukasiewicz many-valued logic. I also provide a new proof of this result, which is purely categorical and does not depend on a particular choice of the dual language. I consider this paper one of my best works as it solves a recent open problem in the field and the proof is elegant and mathematically interesting.
			%%			It is placed in the research area of duality theory. With his Representation Theorem of Boolean Algebras in 1936, M.~H.~Stone linked syntax and semantics by identifying each formula of classical propositional logic with the set of possible worlds in which it stands. Stone’s theorem amounts to a duality (= dual categorical equivalence) between the Boolean algebras and Stone spaces, i.e. compact Hausdorff spaces with a basis of closed open sets. The term “duality” stems from the inversion of the order of morphisms: quotients in the category of algebras (= more information) correspond to subobjects in the category of spaces (= fewer possible worlds), and vice versa.\\			
			%%			As summarized in the Introduction to Johnstone’s book “Stone spaces”, Stone’s paper influenced many areas of mathematics such as logic, functional analysis, topology, category theory, sheaf theory, topos theory, algebraic geometry and representation theory of rings and more general algebraic structures. Moreover, the link between syntax and semantics provided by Stone's duality is particularly central in theoretical computer science since the two sides correspond respectively to specification languages and computational state spaces. The ability to translate faithfully between these two worlds has often proved itself to be a powerful theoretical tool as well as a handle for solving problems. A prime example is Abramsky’s work linking program logic and domain theory via Stone duality. Other examples include large parts of modal and intuitionistic logics, where Jónsson-Tarski duality yields Kripke semantics. Stone duality also played a significant role in automata theory (Pippenger, 1997, Gehrke et al., 2008), computational complexity and the search for lower bounds (Gehrke and Krebs, 2017). Similarly, Abramsky and coworkers connect via Stone duality categorical tools from semantics, such as comonads, with concepts from finite model theory, such as tree-width and tree-depth.\\			
			%			If we drop the assumption of total disconnectedness from Stone spaces, we are left with the category of compact Hausdorff spaces. Duskin showed in 1969 that Stone duality for Boolean algebras can be lifted to compact Hausdorff spaces retaining the algebraic nature. In 1970, H.~A. Priestley generalized Stone's duality to distributive lattices; the dual structures are Priestley spaces, which are certain compact Hausdorff spaces equipped with a partial order relation. A natural generalization of Priestley spaces are compact ordered spaces, introduced by Nachbin; these are the partially ordered version of compact Hausdorff spaces. In a paper by Hofmann, Neves and Nora, the authors left open the problem of determining whether, as with Stone spaces, Priestley duality can be lifted to the category of compact ordered spaces while retaining its algebraic nature. In the paper “The dual of compact ordered spaces is a variety” I showed that this is indeed the case. I did this by providing an explicit equational axiomatization of the dual category, based on an algebraic language related to {\L}ukasiewicz many-valued logic. In this paper with Reggio, we provide a new proof. The proof is purely categorical; thus, it does not depend on a particular choice of the dual language. The proof copes with a difficulty that in the unordered case considered by Duskin is not present (i.e.\  the dual of compact ordered spaces is not Mal’cev), requiring new techniques. 
			%			This elegant and non-trivial proof not only solves a recent open problem, but also makes a mathematically interesting contribution, making this paper one of my best works.
			%%			Thus, besides having an interesting result that solves a recent open problem, the proof is elegant, non-trivial and mathematically interesting; thus I consider this paper one of my best ones.			
			%			\item
			%			Review (MathSciNet):
			%				\href{https://mathscinet.ams.org/mathscinet/pdf/4163309.pdf}{https://mathscinet.ams.org/mathscinet/pdf/4163309.pdf}.
			%		\end{itemize}
	}
	
	
	\cvpub{
		\increaseandshowpubs. M.\ Abbadini.\\
		 Operations that preserve integrability, and truncated {R}iesz spaces.\\
		 {\em \textbf{Forum Mathematicum}}.
		 \gray (2020) 32(6):1487--1513.
		 \gray
		 DOI: \href{https://doi.org/10.1515/forum-2018-0244}{10.1515/forum-2018-0244}.
		%		B\\
		%		\bodyfontlight
		%		\begin{itemize}
			%%			\item 
			%%			Content: For any real number $p \in [1, +\infty)$, we characterize the operations $\mathbb{R}^I \to \mathbb{R}$ that preserve $p$-integrability, i.e.\ those under pointwise application of which the set $L_p(\mu)$ is closed for each measure $\mu$.
			%%			We show that the appropriate variety of algebras whose operations are exactly such functions is the category of $\sigma$-complete Dedekind truncated Riesz spaces.
			%%			Furthermore, we prove that $\mathbb{R}$ generates this variety and we exhibit concrete models of free Dedekind $\sigma$-complete truncated Riesz spaces. 
			%%			Analogous results are obtained for operations that preserve $p$-integrability on finite measure spaces; the corresponding variety coincides with the well-studied category of Dedekind $\sigma$-complete Riesz spaces with weak unit.
			%%%			, $ \mathbb{R}$ is found to generate this variety, and a concrete model of free Dedekind $\sigma$-complete Riesz spaces with weak unit is presented. 
			%%			These results characterize the largest algebraic signature that can be correctly interpreted in all $L_p$ spaces and thus are functional to an algebraic approach to integration.
			%			\item
			%			Review (MathSciNet): \href{https://mathscinet.ams.org/mathscinet/pdf/4168701.pdf}{https://mathscinet.ams.org/mathscinet/pdf/4168701.pdf}.
			%		\end{itemize}
	}
	
	
	
	\cvpub{
		\increaseandshowpubs. M.\ Abbadini.\\
		 Dedekind {$\sigma$}-complete {$\ell$}-groups and {R}iesz spaces as
		varieties.\\
		 {\em \textbf{Positivity}}.
		 \gray (2020) 24(4):1081--1100. 
		 DOI: \href{https://doi.org/10.1007/s11117-019-00718-9}{10.1007/s11117-019-00718-9}.
		%	   B
		%		\bodyfontlight
		%		\begin{itemize}
			%%			\item Content: We prove that the category of Dedekind $\sigma$-complete Riesz spaces is isomorphic to an infinitary variety, we provide an explicit equational axiomatization and we show that $\mathbb{R}$ generates this class as a variety and even as a quasi-variety.
			%%			Analogous results are established for the categories of (i) Dedekind $\sigma$-complete Riesz spaces with a weak order unit, (ii) Dedekind $\sigma$-complete lattice-ordered groups, and (iii) Dedekind $\sigma$-complete lattice-ordered groups with a weak order unit.
			%			\item
			%			Review (MathSciNet): \href{https://mathscinet.ams.org/mathscinet/pdf/4129609.pdf}{https://mathscinet.ams.org/mathscinet/pdf/4129609.pdf}.
			%		\end{itemize}
	}
	
	
	\cvpub{
		\increaseandshowpubs. M.\ Abbadini.\\
		 The dual of compact ordered spaces is a variety.\\
		 {\em \textbf{Theory and Applications of Categories}}.
		 \gray (2019) 34:1401--1439. 
		 URL: \url{http://www.tac.mta.ca/tac/volumes/34/44/34-44abs.html}.%{http://www.tac.mta.ca/tac/volumes/34/44/34-44abs.html}.
		%		B\\
		%		\bodyfontlight
		%		\begin{itemize}
			%%			\item 
			%%			Content: 
			%%			We solve an open problem in 
			%%			``D.~Hofmann, R.~Neves, P.~Nora. Generating the
			%%			algebraic theory of C(X): the case of partially ordered compact spaces. \emph{Theory
				%%				Appl. Categ.}, 33:276–295, 2018''.
			%%%			D.~Hofmann, R.~Neves and P.~Nora proved that the dual of the category of compact ordered spaces is a quasi-variety---non-finitary, but bounded by $\aleph_1$. An open question was: is it also a variety? I show that the answer is yes. 
			%%			We prove that the dual of the category of compact ordered spaces is a quasi-variety---not finitary, but bounded by $\aleph_1$.
			%%			We describe the variety via a set of finitary operations with an operation of countably infinite arity and equational axioms.
			%			\item
			%			Review (MathSciNet): \href{https://mathscinet.ams.org/mathscinet/pdf/4041843.pdf}{https://mathscinet.ams.org/mathscinet/pdf/4041843.pdf}.\\
			%			Review (Zentral Blatt): \href{https://zbmath.org/?q=an\%3A1446.18003}{https://zbmath.org/?q=an\%3A1446.18003}.
			%		\end{itemize}		
	}
\end{cvpubs}


\medskip



%---------------------------------------------------------
\cvsubsection{\stress{Peer-reviewed book chapters}}
%---------------------------------------------------------



\begin{cvpubs}
	\small 
	
	
	\cvpub{
		\increaseandshowpubs. M.\ Abbadini, D.\ Mundici.\\
		 Unital Specker $\ell$-groups and boolean multispaces.\\
		 In \emph{\textbf{volume dedicated to Hilary Priestley}} (title to be announced).
		 \gray Outstanding Contributions to Logic, Springer (to appear). Edited by B. Davey, L. M. Cabrer, A. Palmigiano.
		 Preprint at \href{https://arxiv.org/abs/2508.21500}{arXiv:2508.21500}.
	}
	
	
\end{cvpubs}


\medskip

% %---------------------------------------------------------
\cvsubsection{\stress{Proceedings in international peer-reviewed conferences}}
% %---------------------------------------------------------

\begin{cvpubs}
	\small 
	
	
	\cvpub{
		\increaseandshowpubs. M.\ Abbadini, F.\ Di~Stefano, L.\ Spada.\\
		 Unification in {\L}ukasiewicz logic with a finite number of
		variables.\\
		 In {\em \textbf{Information Processing and Management of Uncertainty in
				Knowledge-Based Systems}}.
				\gray (2020) 1239:622--633. Springer International Publishing. 
		 DOI: \href{https://doi.org/10.1007/978-3-030-50153-2_46}{10.1007/978-3-030-50153-2\_46}.
%		\bodyfontlight
		%\begin{itemize}
		%	\item Content: The study of unification modulo an equational theory has gained increasing importance in recent years. The unification problem for the logic of {\L}ukasiewicz is: given two terms $s$, $t$, find a uniform replacement (called a unifier) of the variables appearing in $s$ and $t$ by terms that makes $s$ and $t$ identical modulo the theory of {\L}ukasiewicz logic.
		%	The most fundamental information one wants is whether two unifiable terms admit a most general unifier.
		%	We prove that the logic of {\L}ukasiewicz with a finite number of variables does not have this property. 
		%	More precisely, we prove that the type of unification of Lukasiewicz logic with a finite number of variables is neither unary nor finitary.
		%\end{itemize}
	}
	
	
\end{cvpubs}


\medskip

% %---------------------------------------------------------
\cvsubsection{\stress{Preprints}}
% %---------------------------------------------------------

\begin{cvpubs}

	\cvpub{
		\increaseandshowpubs. M.\ Abbadini, A.\ Jung.\\
		 On the symmetry behind duality.\\
		 \gray Under review. Preprint at \href{https://arxiv.org/abs/2507.18245}{arXiv:2507.18245}.
	}

	
	\cvpub{
		\increaseandshowpubs. M.\ Abbadini, A.\ P\v{r}enosil.\\
		 Duality for finitely valued algebras.\\
		 \gray Under review. Preprint at \href{https://arxiv.org/abs/2505.11490}{arXiv:2505.11490}.
	}
	
		

	\cvpub{
		\increaseandshowpubs. M.\ Abbadini, F.\ Guffanti.\\
		 Freely adding one layer of quantifiers to a Boolean doctrine.\\
		 \gray Under review. Preprint at \href{https://arxiv.org/abs/2410.16328}{arXiv:2410.16328}.
	}

	
\end{cvpubs}
